\section{The Single Audit as a history-dependent observation regime}\label{sec:regime}

The empirical illustration comes from the Single Audit, and its terms need
defining before the data are read.

A Single Audit is an organization-wide audit required of nonfederal entities,
such as states, local governments, universities and nonprofits, that expend
federal awards above a regulatory threshold. It combines a financial statement
audit with compliance testing of the entity's federal programs, and its
structured results are published by the Federal Audit Clearinghouse.

Compliance testing is not spread evenly across every program. The auditor
designates certain federal programs as \emph{major programs} and performs the
required major-program compliance audit on them. Programs not selected as
major do not receive that full program audit solely by virtue of program
status, although other work, including required follow-up of prior findings,
may still occur. The selection is risk based. Programs whose federal
expenditures exceed the Type~A threshold are \emph{Type~A}; the rest are
Type~B. The threshold is tiered by the entity's total federal expenditures:
over the sample period it was \$750,000 for entities expending up to \$25
million and higher in tiers above that (2~CFR 200.518(b)(1)). A Type~A program that meets a list of low-risk criteria may be
left out of detailed testing in a given year, and one that fails the criteria
must be tested as major. Prior findings are among
the criteria: a material weakness, a modified opinion on the program, or
questioned costs above five percent of the program's expenditures in the most
recent audit period disqualify a Type~A program from low-risk
treatment.\footnote{Not every prior finding does this. A significant
deficiency that is not a material weakness, or questioned costs below the
five-percent line, leaves a Type~A program eligible for low-risk treatment.
The paper does not assume that every prior finding makes major-program
coverage more likely, and its sample construction conditions on major-program
coverage in both years.} The
requirements tested for a major program are those the Office of Management and
Budget's annual \emph{Compliance Supplement} designates as applicable to it.

Separately, and regardless of program selection, the auditor must follow up on
prior audit findings, whether or not the program with the prior finding is
selected as major again. In the governing regulation, prior findings raise
assessed program risk (2~CFR 200.519(b)(2)); they can disqualify a Type~A
program from low-risk treatment (200.518(c)(1)); and follow-up is mandatory
``regardless of whether a prior audit finding relates to a major program in
the current year'' (200.514(e)). None of this is a defect; the follow-up requirement exists
because findings left uncorrected are a documented problem in this program
\citep{gao2002}.

The setting is useful for one reason. In most inspection systems the feedback
from past results to present attention has to be inferred from inspector
behavior. Here it is written into the rules, through two channels: a coverage
channel, in which a prior finding can change whether a program is tested in
depth at all, and a condition-specific channel, in which a prior finding must
be followed up wherever it sits. Restricting the sample to programs audited as
major in both years conditions on the major-program-selection channel: every
comparison is among programs selected in both years. That conditioning does
not make the selection ignorable, and it does not hold audit intensity or
condition-specific observation fixed: the
follow-up channel is mandated inside the restricted sample, so two such
programs, one with a prior finding and one without, face different
condition-specific attention by regulation even though both are audited in
depth in both years.
